\documentclass[11pt,a4paper]{article}
\usepackage[margin=25mm]{geometry}
\usepackage{amsmath,amssymb,amsthm,hyperref}
\setlength{\emergencystretch}{2em}
\newtheorem{theorem}{Theorem}
\newtheorem{lemma}[theorem]{Lemma}
\newcommand{\gen}[1]{\langle#1\rangle}
\title{Centralizers of a maximal block action}
\author{Human Compression project: a component of the L044 recovery}
\date{}
\begin{document}\maketitle
Let $G=A_n=\gen{\rho_0,\ldots,\rho_{r-1}}$ in its natural action,
with the distinguished involutions satisfying the string commutations and
intersection property. Put $H=G_i=\gen{\rho_j:j\ne i}$.
This paper proves the following part of the main rank bound.

\begin{theorem}\label{main}
Suppose $n\ge12$, $H$ is transitive, and a maximal nontrivial block
system for $H$ has $m>2$ blocks of size $k>2$. Then
\[
                         r\le\lfloor(n-1)/2\rfloor.
\]
\end{theorem}

The source is Cameron--Fernandes--Leemans--Mixer, \emph{Highest rank
of a polytope for $A_n$}, arXiv:1605.09173v1, Proposition~3.2.
We acquire the complete maximal-block argument from the reviewed L047
recovery, including Whiston's independent-family bound and its equality
clause. We also acquire two precisely bounded finite conclusions: no
transitive subgroup of $A_6$ has a rank-four string C-marking, from the
complete L047 proper-group census and its separate $A_6$ exclusion;
and $A_{12}$ has no marking of rank greater than five, from the complete
L071 census. These are retained theorems with their proofs and execution
records, not new computations here.

The new work below is an equality argument for the block decomposition.
In particular, we do not require the general assertion that an omitted
mark cannot commute with a subgroup merely because its block image is
primitive. Faithfulness and the centralizer are established when needed.

\section{The block decomposition and its equality cases}
Choose an inclusion-minimal distinguished subfamily $L$ whose images
generate the primitive block group $P$. Let $C$ be the distinguished
marks outside $L$ commuting with all of $L$, and let $R$ be the remaining
marks of $H$. Write $B=\gen L$ and $D=\gen C$.
The acquired argument has the following exact consequences.
If the image diagram of $L$ is disconnected, then $m\ge60$ and
\[
                    r\le2\log_2m+k+4.
\]
This is smaller than $(km-1)/2$ for $k\ge3,m\ge60$: the difference
is increasing in both variables and is already positive at $(3,60)$.
We may therefore suppose that $L$ occupies an interval of labels.
In that case
\[
 |L|\le m-1,\qquad |C|\le k-1,\qquad |R|\le2,
 \qquad r\le k+m.                                      \tag{1}
\]
The last inequality includes the acquired exclusion of
$\operatorname{rank}(H)=k+m$.

The group $D$ fixes each block and has a faithful restriction to every
block. Indeed, the primitive group $P$, generated by involutions, has
trivial centralizer in $S_m$ when $m>2$. A nontrivial primitive centralizer
would force $P$ to be regular of prime degree; the cyclic group of odd
prime order is not generated by involutions. Thus $D$ fixes the blocks.
An element of $D$ fixing one block pointwise fixes every block, by
commutation with the block-transitive group $B$. This proves faithfulness.
Transporting the coordinates of one block by elements of $B$ identifies
$D$ with the same permutation group on each block.

The intersection property gives $B\cap D=1$, and every distinguished
subfamily remains independent. If $|C|=k-1$, Whiston's equality clause
therefore makes $D$ the diagonal $S_k$. Its centralizer in the full
block stabilizer is $S_m$ acting consistently on the block coordinates.
Consequently $B$ is faithful on the blocks and is exactly the diagonal
copy of $P$ in these coordinates.

\begin{lemma}\label{equality}
If $r=k+m$, then
\[
                  |L|=m-1,\qquad |C|=k-2,\qquad |R|=2.   \tag{2}
\]
In particular $P=S_m$ and $\rho_i$ commutes with $B$.
\end{lemma}
\begin{proof}
Equation~(1) and $r-1=|L|+|C|+|R|$ leave exactly three possibilities:
\[
 (|L|,|C|,|R|)=(m-1,k-1,1),\ (m-2,k-1,2),\ (m-1,k-2,2).
\]
Consider the second. Here $B$ is a faithful diagonal copy of the primitive
group $P$. Since $P$ has trivial centralizer, the full centralizer of $B$
on the $km$ points consists of permutations of its $k$ identical rows,
acting consistently on them. It fixes each of the original blocks.
Both neighbours of the interval $L$ are occupied by the two marks in $R$.
Hence the omitted label $i$ is not adjacent to $L$, so $\rho_i$ centralizes
$B$ and preserves the blocks. Then $G=\gen{H,\rho_i}$ preserves them,
contrary to primitivity.

In the first possibility, $B=S_m$ and $D=S_k$ form their product action.
The noncommuting diagram of $C$ is connected: disconnected components
would give a nontrivial direct product decomposition of the centerless,
directly indecomposable group $S_k$. Thus $C$, as well as $L$, is an
interval of labels. The one mark in $R$ is adjacent to $L$. It cannot
centralize $D$, since its membership in the block stabilizer would then
put it in $B$, contradicting independence. It is therefore also adjacent
to $C$. The two intervals and this one intervening mark occupy all the
$r-1$ labels of $H$ consecutively. The omitted label is at an end and
centralizes at least one of $B,D$ in full. A permutation centralizing
$D$ preserves its orbits, which are the original blocks. A permutation
centralizing the diagonal $S_m$ fixes each original block, as in the
previous paragraph. Either way $G$ preserves the block system, again a
contradiction. Only~(2) remains. Its last two conclusions follow from
Whiston's equality clause and the two occupied neighbours of $L$.
\end{proof}

\section{Three elementary centralizer arguments}
\begin{lemma}\label{constants}
Let $V$ be a vector space over $\mathbb F_2$ and let $m$ be even.
Every nonzero $S_m$-invariant subspace of $V^m$, with $S_m$ permuting
coordinates, contains a nonzero constant vector.
\end{lemma}
\begin{proof}
A nonzero constant vector already has the required form. Otherwise take
a vector whose $i$ and $j$ coordinates differ. Adding its image under
$(i,j)$ gives a vector with the same nonzero value $v$ at $i,j$ and zero
elsewhere. Its conjugates give $v$ at any chosen pair of positions.
Partition the $m$ positions into pairs and add the corresponding vectors.
The result is $(v,\ldots,v)$.
\end{proof}

\begin{lemma}\label{three}
Let $m\ge4$ be even. Suppose $B\le A_{3m}$ preserves $m$ blocks of
size three and induces $S_m$ on them. Let $z$ be a permutation acting
as a transposition in each block. If $B$ centralizes $z$, then $z\in B$.
\end{lemma}
\begin{proof}
Choose coordinates in which $z$ is the same transposition on each block.
The centralizer of that transposition in $S_3$ is its own cyclic group.
Thus $B$ lies in $C_2^m\rtimes S_m$, in the degree-$3m$ action with one
fixed point per block for the base group. Its kernel $K$ on blocks is an
$S_m$-invariant subspace of $\mathbb F_2^m$. If $K\ne0$, Lemma~\ref{constants}
puts its nonzero constant vector $z$ in $K$, as claimed.

If $K=0$, the projection identifies $B$ with $S_m$. Let $t\in B$ be
the unique lift of a transposition of two blocks. It is an involution,
and on the two exchanged blocks it is a product of three transpositions.
The lifts of the permutations of the other $m-2$ blocks centralize $t$
and act transitively on those blocks. Consequently its restrictions to
the fixed blocks are conjugate. Each restriction is either the identity
or the transposition defining $z$, so all have the same parity. There
are an even number $m-2$ of these blocks. Their total contribution is
even, whereas the contribution from the exchanged blocks is odd.
Thus $t$ is odd, contradicting $B\le A_{3m}$.
\end{proof}

\begin{lemma}\label{fourfaithful}
Suppose $B$ acts on sixteen points, preserves four blocks of size four,
and maps isomorphically onto $S_4$ in its block action. If a nontrivial
group $D$ fixes the blocks and centralizes $B$, then the full centralizer
of $B$ fixes each block.
\end{lemma}
\begin{proof}
Identify $B$ with $S_4$ through the block action and fix one block.
Its stabilizer is $S_3$. A $B$-orbit meets this block in an orbit of that
$S_3$, of length $d\le4$ dividing six. Thus $d=1,2,3$ and the full
orbit has length $4d$. The corresponding point stabilizers are
$S_3,A_3,C_2$, respectively.

If a twelve-point orbit occurs, the remaining orbit has four points.
On the twelve-point orbit, permutations commuting with $B$ and fixing
blocks correspond to $N_{S_3}(C_2)/C_2=1$. On the four-point orbit
the analogous quotient is $N_{S_3}(S_3)/S_3=1$. Unequal orbit sizes
prevent any interchange of the two orbits. This would make $D$ trivial,
a contradiction.

All orbits therefore have length four or eight. For their point
stabilizers $J=S_3,A_3$ we have $N_{S_4}(J)=S_3$, the block stabilizer.
In the coset action, a commuting permutation is given by an element of
$N_{S_4}(J)/J$ and hence fixes every block. The same calculation applies
to equivariant bijections between two isomorphic orbits: choose their
base points in the same original block, so both stabilizers are the
same $J\le S_3$. The element defining an equivariant bijection again
belongs to $N_{S_4}(J)\le S_3$. Every permutation centralizing $B$ is
assembled from these within-orbit maps and equivariant interchanges.
It therefore fixes each block, as asserted.
\end{proof}

\begin{lemma}\label{fourkernel}
Suppose $B,D$ preserve four blocks of size four, commute, and have
$B\cap D=1$. Suppose $B$ induces $S_4$ on the blocks, while $D$ fixes
them and has a faithful restriction generated by two independent
involutions. Then $B$ acts faithfully on the blocks, and its full
centralizer fixes each block.
\end{lemma}
\begin{proof}
Transport coordinates by $B$ so that $D$ acts identically on all blocks.
Two distinct independent involutions in $S_4$ generate a dihedral group
of order four, six or eight, since the order of their product is two,
three or four. The order-six group is an $S_3$ fixing a point, and has
trivial centralizer in $S_4$. The order-eight group is a Sylow subgroup,
whose centralizer is its central involution. The order-four groups are
either the regular Klein group or the group generated by two disjoint
transpositions; their centralizer is the group itself. These elementary
centralizers can be obtained by preserving the orbits and commuting with
the two generators. In every case $Z=C_{S_4}(D)$ is an elementary
abelian two-group contained in $D$, possibly trivial.

Now $B\le Z^4\rtimes S_4$. Since $Z$ is abelian, conjugation by $B$
on its block kernel $K\le Z^4$ is the ordinary permutation of coordinates.
If $K\ne0$, Lemma~\ref{constants} supplies a nontrivial diagonal
element of $Z$ in $K$. That element belongs to both $B$ and $D$, a
contradiction. Thus $K=0$. Lemma~\ref{fourfaithful} now applies.
\end{proof}

\section{Small even actions}
\begin{lemma}\label{five}
Three independent involutions in $A_5$ generate $A_5$.
\end{lemma}
\begin{proof}
An intransitive subgroup is contained in $A_4$ or in the even part of
$S_3\times S_2$, which is $S_3$. In $A_4$ all involutions lie in its
Klein subgroup; neither this group nor $S_3$ has three independent
involutions. A proper transitive subgroup has order divisible by five
and dividing sixty, hence order $5,10,15,20$ or $30$. Orders twenty
and thirty would give faithful nontrivial actions of the simple group
$A_5$ on three or two cosets, which is impossible. A group of order
fifteen has normal Sylow three- and five-subgroups and hence an element
of order fifteen, absent from $A_5$. Order five supplies no involutions;
at order ten, two distinct involutions generate the whole dihedral group.
The only remaining possibility is $A_5$ itself.
\end{proof}

\begin{lemma}\label{six}
No subgroup of $A_6$ has a string C-marking of rank four.
\end{lemma}
\begin{proof}
The transitive case is precisely the acquired L047 finite conclusion
specified above. For an intransitive group choose a nonempty invariant
set of size $a\le3$, passing to its complement if necessary. It is
contained in $(S_a\times S_{6-a})\cap A_6$. For $a=1$ this is $A_5$,
whose independent length is at most three by Whiston's equality clause
in $S_5$. For $a=2$ it is isomorphic to $S_4$, since the permutation
on the two-point factor is determined by the sign on the other factor;
Whiston again bounds the independent length by three. For $a=3$ its
order is eighteen, so every strictly increasing subgroup chain has at
most three steps. An independent family produces such a chain and has
at most three members. All cases exclude rank four.
\end{proof}

\section{Proof of the rank bound}
\begin{proof}[Proof of Theorem~\ref{main}]
At $n=12$ the complete acquired $A_{12}$ census gives $r\le5$.
Assume $n>12$. The disconnected case was settled before~(1).
If~(1) does not already imply $r\le\lfloor(km-1)/2\rfloor$, then
\[
        (k,m)\in\{(3,5),(3,6),(4,4),(5,3),(6,3)\}.
\]
Indeed, $2(k+m)\ge km$ is necessary when $km$ is even; when it is
odd the failed integer bound requires $2(k+m)\ge km+1$.
In particular $(k-2)(m-2)\le4$, and listing integers $k,m>2$ with
$km>12$ gives the five pairs above. In each case failure forces
$r=k+m$. Lemma~\ref{equality} applies, so $P=S_m$, $|C|=k-2$,
$|R|=2$, and $\rho_i$ centralizes $B$.

For $(k,m)=(3,5)$, $D$ has one involutory generator. Its faithful
restriction to a block is a transposition. It acts as such on all five
blocks and is therefore odd, a contradiction.

For $(3,6)$, write that generator as $z$. Lemma~\ref{three} puts
$z$ in $B$, contrary to $B\cap D=1$.

For $(5,3)$, the three independent involutions of $D$ restrict to
even permutations on a block, since the number of identical blocks is
odd. Lemma~\ref{five} gives $D=A_5$ diagonally. Its centralizer in
$S_5$ is trivial, so $B$ is the diagonal faithful $S_3$ on the three
blocks. A block transposition is a product of five transpositions on
points and is odd. Again this contradicts $B\le A_{15}$.

For $(6,3)$, the four generators of $D$ similarly restrict to a
string C-marking of rank four in $A_6$. Restriction is faithful and
preserves the intersection property. Lemma~\ref{six} excludes it.

Finally, for $(4,4)$, $D$ has two independent involutory generators.
Lemma~\ref{fourkernel} says that every permutation centralizing $B$
fixes each original block. In particular $\rho_i$ does so. Together
with $H$, it generates a group preserving that nontrivial block system,
contradicting $G=A_{16}$. This proves the theorem in every case.
\end{proof}

\paragraph{Scope and attribution.}
The source's rank conclusion is recovered only when both block parameters
exceed two. The cases of two-point blocks and of two blocks remain
separate obligations. This component uses no induction on alternating
degree and no new finite search. The L047 decomposition, Whiston input,
classical simplicity of $A_5$, and acquired $A_6,A_{12}$ classifications
remain explicitly borrowed. The equality and centralizer arguments above
are this continuation; no priority claim is made for their elementary
ingredients. No whole-original, open-question, Lean or numerical
compression credit is assigned by this component.
\end{document}
